\section{Discussion}
\label{sec:discussion}

\para{What is the strongest claim supported by the data?} The strongest claim is that L2 regularization reliably reduces overfitting in this small-data setting. The average train-validation gap falls by about $0.104$, the paired confidence interval excludes zero, and the effect appears in $11/12$ runs. That is exactly the pattern one would expect if shrinkage were acting mainly as variance control.

\para{Why is the accuracy effect weaker?} Accuracy combines both variance reduction and bias increase. On easier datasets such as \breastcancer and \wine, shrinkage reduces variance without paying much cost in bias, so both predictive performance and gap metrics improve. On more difficult datasets such as \sonar and \glass, the same shrinkage can remove enough flexibility to offset any variance benefit. The result is a smaller gap but no reliable accuracy gain. This interpretation is aligned with the conditional view of regularization in prior work \cite{zhang2017understanding,power2022grokking,dangelo2024weight}.

\para{What does this imply for practice?} For small tabular datasets, L2 is well justified when the goal is to stabilize a linear model and control apparent overfitting. It is less justified if the claim is that it will reliably increase held-out accuracy. Practitioners should therefore treat L2 as a dependable shrinkage tool and tune it carefully, rather than assuming it is a universal improvement over a weakly regularized baseline.

\para{Limitations.} This study has four important limitations. First, it only evaluates logistic regression, so the results may not transfer directly to multilayer perceptrons or tree-based models. Second, the aggregate analysis covers only 12 paired observations, and the per-dataset analyses use only three seeds each. Third, the baseline is only approximately unregularized because it uses \texttt{C=1e12} for compatibility with the current \texttt{scikit-learn} interface. Fourth, the split protocol uses a single validation split per seed rather than nested cross-validation, so selection variance may still be understated. In addition, the Shapiro--Wilk test rejects normality for the aggregate train-gap differences, which suggests caution when reading the corresponding paired $t$-test even though the effect size is large.

\para{Broader implications.} The broader lesson is methodological. In small-sample settings, it is easy to overstate a regularization result if one reports only test accuracy on one split. Our results show why paired repeated-split analysis matters: the same intervention can have a strong and statistically supported effect on overfitting indicators while failing to deliver a statistically supported gain on final accuracy.
